\documentclass[11pt]{article}
\usepackage{mathblatt}
\begin{document}
\weit
\typeout{PROBE-BEGIN flaechen-e3-k1-s0-v1}
\begin{aufgabe}{\detokenize{flaechen-e3-k1-s0-v1}}
\begin{teile}
\teil Probe

\dreieck{(0,0)}{(5,0)}{(2,3)}{}{}{$g$}{}{}{}
\end{teile}
\end{aufgabe}
\typeout{PROBE-END flaechen-e3-k1-s0-v1}
\typeout{PROBE-BEGIN pythagoras-e1-k1-s0-v2}
\begin{aufgabe}{\detokenize{pythagoras-e1-k1-s0-v2}}
\begin{teile}
\teil Probe

\dreieck{(0,0)}{(4.5,0)}{(1.5,2.8)}{}{}{}{$62^\circ$}{$48^\circ$}{$70^\circ$}
\end{teile}
\end{aufgabe}
\typeout{PROBE-END pythagoras-e1-k1-s0-v2}
\typeout{PROBE-BEGIN symmetrie-abbildungen-e2-k2-s0-v4}
\begin{aufgabe}{\detokenize{symmetrie-abbildungen-e2-k2-s0-v4}}
\begin{teile}
\teil Probe

\dreieck{(0,0)}{(4,0)}{(2,3)}{a}{b}{c}{\alpha}{\beta}{\gamma}
\end{teile}
\end{aufgabe}
\typeout{PROBE-END symmetrie-abbildungen-e2-k2-s0-v4}
\typeout{PROBE-BEGIN terme-e1-k1-s3-v2}
\begin{aufgabe}{\detokenize{terme-e1-k1-s3-v2}}
\begin{teile}
\teil Probe

\dreieck{(0,0)}{(6,0)}{(3,4)}{x}{x}{6}{}{}{}
\end{teile}
\end{aufgabe}
\typeout{PROBE-END terme-e1-k1-s3-v2}
\typeout{PROBE-BEGIN trigonometrie-e1-k1-s0-v1}
\begin{aufgabe}{\detokenize{trigonometrie-e1-k1-s0-v1}}
\begin{teile}
\teil Probe

\dreieck{(0,0)}{(5,0)}{(5,3)}{a}{?}{c}{\alpha}{}{}
\end{teile}
\end{aufgabe}
\typeout{PROBE-END trigonometrie-e1-k1-s0-v1}
\end{document}
